%% ============================================================
%% Chaos Agent v2.2 — Technical Documentation
%% Author: Ivan Putna (Architect of Chaos — AoCH)
%% Compiler: XeLaTeX (Overleaf default)
%% ============================================================

\documentclass[11pt, a4paper]{article}

%% ── Fonts & encoding ────────────────────────────────────────
\usepackage{fontspec}
\setmainfont{EB Garamond}
\setsansfont{Fira Sans}[BoldFont = Fira Sans SemiBold]
\setmonofont{Fira Code}[Scale = 0.88]

%% ── Page geometry ───────────────────────────────────────────
\usepackage[
  a4paper,
  top=28mm, bottom=28mm,
  left=32mm, right=26mm,
  marginparwidth=18mm
]{geometry}

%% ── Colours ─────────────────────────────────────────────────
\usepackage{xcolor}
\definecolor{accent}{HTML}{1A1A2E}
\definecolor{rule}{HTML}{C8A96E}
\definecolor{codeframe}{HTML}{EDEDED}
\definecolor{linkcolor}{HTML}{2E4057}
\definecolor{note}{HTML}{6B6B6B}
\definecolor{added}{HTML}{2D6A4F}
\definecolor{changed}{HTML}{7B5EA7}

%% ── Hyperlinks ──────────────────────────────────────────────
\usepackage[colorlinks=true, linkcolor=linkcolor, urlcolor=linkcolor, citecolor=linkcolor]{hyperref}

%% ── Typography ──────────────────────────────────────────────
\usepackage{microtype}
\usepackage{setspace}
\setstretch{1.22}
\usepackage{parskip}
\setlength{\parskip}{6pt}

%% ── Section styling ─────────────────────────────────────────
\usepackage{titlesec}
\titleformat{\section}
  {\sffamily\large\bfseries\color{accent}}{}{0em}{}
  [\vspace{1pt}{\color{rule}\rule{\linewidth}{0.6pt}}]
\titleformat{\subsection}
  {\sffamily\normalsize\bfseries\color{accent}}{}{0em}{}
\titlespacing{\section}{0pt}{18pt}{6pt}
\titlespacing{\subsection}{0pt}{12pt}{4pt}

%% ── Table of contents ───────────────────────────────────────
\usepackage{tocloft}
\renewcommand{\cftsecfont}{\sffamily\small}
\renewcommand{\cftsecpagefont}{\sffamily\small\color{note}}
\renewcommand{\cftdotsep}{4}
\setlength{\cftbeforesecskip}{3pt}

%% ── Code listings ───────────────────────────────────────────
\usepackage{listings}
\lstdefinestyle{chaos}{
  language=Python,
  backgroundcolor=\color{codeframe},
  basicstyle=\ttfamily\footnotesize,
  keywordstyle=\bfseries\color{accent},
  commentstyle=\itshape\color{note},
  stringstyle=\color{linkcolor},
  numberstyle=\tiny\color{note},
  numbers=left, numbersep=8pt,
  frame=none, framesep=6pt,
  xleftmargin=14pt,
  breaklines=true,
  showstringspaces=false,
  tabsize=4, captionpos=b,
  aboveskip=10pt, belowskip=6pt,
}
\lstset{style=chaos}

%% ── Tables ──────────────────────────────────────────────────
\usepackage{booktabs}
\usepackage{tabularx}
\usepackage{array}
\renewcommand{\arraystretch}{1.3}

%% ── Margin notes ────────────────────────────────────────────
\usepackage{marginnote}
\renewcommand{\marginfont}{\footnotesize\itshape\color{note}}

%% ── Utilities ───────────────────────────────────────────────
\usepackage{enumitem}
\setlist[itemize]{leftmargin=*, itemsep=2pt, topsep=4pt}
\setlist[enumerate]{leftmargin=*, itemsep=2pt, topsep=4pt}
\usepackage{graphicx}

%% ── Header / footer ─────────────────────────────────────────
\usepackage{fancyhdr}
\pagestyle{fancy}
\fancyhf{}
\fancyhead[L]{\sffamily\footnotesize\color{note} Chaos Agent v2.2}
\fancyhead[R]{\sffamily\footnotesize\color{note} Ivan Putna — AoCH}
\fancyfoot[C]{\sffamily\footnotesize\color{note} \thepage}
\renewcommand{\headrulewidth}{0pt}

%% ============================================================
\begin{document}
\thispagestyle{empty}

%% ── Title block ─────────────────────────────────────────────
\vspace*{18mm}
{\color{rule}\rule{\linewidth}{1.2pt}}
\vspace{6pt}

{\sffamily\Huge\bfseries\color{accent} Chaos Agent v2.2}

\vspace{4pt}
{\sffamily\large\color{note} Technical Architecture, Design Rationale \& Changelog}

\vspace{6pt}
{\color{rule}\rule{\linewidth}{0.4pt}}
\vspace{10mm}

\begin{tabular}{@{}ll}
  {\sffamily\small\color{note} Author}    & Ivan Putna \\
  {\sffamily\small\color{note} Pseudonym} & Architect of Chaos (AoCH) \\
  {\sffamily\small\color{note} Version}   & 2.2 \\
  {\sffamily\small\color{note} v2.0 date} & 2026-03-31 \\
  {\sffamily\small\color{note} v2.1 date} & 2026-04-01 \\
  {\sffamily\small\color{note} v2.2 date} & 2026-04-01 \\
  {\sffamily\small\color{note} License}   & MIT \\
  {\sffamily\small\color{note} GitHub}    &
    \href{https://github.com/iVANIUS1}{github.com/iVANIUS1} \\
  {\sffamily\small\color{note} LinkedIn}  &
    \href{https://www.linkedin.com/in/ivan-putna-3313a3381}{linkedin.com/in/ivan-putna-3313a3381} \\
\end{tabular}

\vspace{14mm}
\tableofcontents
\newpage

%% ============================================================
\section{Philosophy}

\begin{quote}
  \itshape Less code. More intelligence.
\end{quote}

Chaos Agent is a minimal, production-ready AI coding agent built around a single
conviction: complexity is not robustness. The architecture demonstrates that a
system with one source of truth, an explicit cognitive layer, and graceful
degradation can be more reliable than a heavily layered alternative — while
remaining dramatically easier to reason about.

Each version is documented with the rationale behind every decision. This is not
a framework. It is a reference architecture — designed to be read, understood,
and adapted.

%% ============================================================
\section{Architecture Overview}

\begin{center}
\begin{tabular}{@{} >{\sffamily\bfseries}l p{9.5cm} @{}}
  \toprule
  Component        & Responsibility \\
  \midrule
  ChaosKernel      & Single source of truth. Owns \texttt{emit()}, \texttt{TokenBudget},
                     compression, and checkpointing. \\
  GuardianAgent    & Centralised policy enforcement via \texttt{GuardianPolicy} and
                     \texttt{audit\_code()}. \\
  ResearcherAgent  & Technical research phase before any implementation. \\
  EngineerAgent    & Code generation, informed by research output. \\
  BaseAgent        & Zero-boilerplate shared \texttt{emit()} API for all agents. \\
  \texttt{audit\_code()}   & Standalone single-pass AST safety auditor. \\
  \texttt{export\_bundle()} & SHA-256 reproducibility bundle (ZIP + manifest). \\
  \texttt{verify\_bundle()} & Integrity verification against manifest. \\
  \bottomrule
\end{tabular}
\end{center}

\subsection{Event data flow}

\begin{enumerate}
  \item An agent or human caller invokes \texttt{kernel.emit(event)}.
  \item The kernel enforces the rate limit, then stamps metadata
        (id, timestamp, depth, active signals).
  \item The event is appended to \texttt{kernel.history} — the canonical log.
  \item The output handler is called (built-in filter or custom callback).
  \item If depth \emph{or} token budget pressure exceeds threshold, compression fires.
  \item Every $N$ events, a checkpoint is written atomically to disk.
\end{enumerate}

There are no side channels. There is no other way for state to enter the system.

%% ============================================================
\section{Core Design Decisions — v2.0}

\subsection{Single \texttt{emit()} gating point}

All events, logs, and state mutations flow through one function. Debugging
requires tracing one call site. Side effects are explicit.

\begin{lstlisting}[caption={All events enter through a single function}]
async def emit(self, event: Dict[str, Any]) -> None:
    # rate limit, stamp metadata, append to history,
    # call output handler, trigger compression if needed
    ...
\end{lstlisting}

\marginnote{Single responsibility applied at the system level, not just the class level.}

Every run produces a complete \texttt{Artifact} — a serialisable snapshot of
every event that occurred, including signals, depth, and token count.

\subsection{Explicit cognitive signals}

Purely reactive agents have no mechanism for intent injection at runtime.
Chaos Agent treats cognitive signals as first-class events:

\begin{lstlisting}[caption={Signals shape reasoning before execution begins}]
await kernel.inject("DOUBT",
    "Direct recursion without memoization will explode on large n.")
await kernel.inject("DIRECTION",
    "Prioritise clarity, safety, and minimalism.")
\end{lstlisting}

Signals accumulate in \texttt{kernel.signals} and are included in every
subsequent event's metadata. They provide active context without polluting the
main message history.

\subsection{Graceful anti-entropy}

Hard depth limits abort long-running tasks at the worst possible moment.
Chaos Agent detects cognitive overload and responds with compression.

\begin{lstlisting}[caption={Depth resets gracefully — no exception raised}]
async def compress(self) -> None:
    # Summarise recent history via LLM call
    # Replace tail of history with summary
    # Reset depth: max(3, depth - 5)
    # Net per cycle: -5 + 2 internal emits = -3 (stable)
    ...
\end{lstlisting}

The \texttt{self.compressed} flag prevents re-entrant compression during the
compression call itself.

\subsection{Declarative \texttt{GuardianPolicy}}

\begin{lstlisting}[caption={Policy is data — swap without touching agent logic}]
@dataclass
class GuardianPolicy:
    allowed_commands: List[str]
    allowed_paths: List[str]
    max_file_size_mb: int = 10
    no_shell: bool = True
    forbid_risky_imports: bool = False
\end{lstlisting}

Every action passes through \texttt{GuardianAgent.check()} before execution.

\subsection{Pluggable output strategy}

\begin{lstlisting}[caption={Two strategies coexist — caller chooses}]
# Built-in attention filter (default)
kernel = ChaosKernel()

# Full control via custom async handler
kernel = ChaosKernel(on_event=my_async_handler)
\end{lstlisting}

\texttt{\_should\_show\_to\_user()} is a method on the kernel — overridable or
replaceable entirely by supplying \texttt{on\_event}. This mirrors Python's
\texttt{logging} design: sensible defaults, no ceiling on customisation.

%% ============================================================
\section{v2.1 Additions}

\subsection{AST-level safety audit}

\marginnote{\color{changed} Changed from v2.0}

String-matching (v2.0) can be bypassed via concatenation or aliasing.
\texttt{audit\_code()} runs a single-pass AST visitor that resolves import aliases
and detects obfuscation:

\begin{lstlisting}[caption={Detects what string-matching cannot}]
# Blocked: exec/eval/compile
exec("rm -rf /")

# Blocked: aliased dangerous attribute
import os as o; o.system("ls")

# Blocked: string obfuscation
m = 'o' + 's'    # combined == "os" -> hard violation

# Blocked: shell=True (when no_shell=True)
subprocess.run(["ls"], shell=True)

# Allowed: benign attribute
x = getattr(obj, "join")
\end{lstlisting}

\begin{lstlisting}[caption={Public API for the auditor}]
allowed, findings, meta = audit_code(
    source,
    no_shell=True,             # block shell=True at AST level
    forbid_risky_imports=False # default off: import os is allowed,
                               # calling os.system is not
)
\end{lstlisting}

\subsection{EMA-based \texttt{TokenBudget}}

\marginnote{\color{added} New in v2.1}

Depth-only compression triggers fire on isolated spikes. \texttt{TokenBudget}
tracks token usage with an exponential moving average (EMA, $\alpha = 0.35$)
and fires compression only on sustained pressure:

\[
\text{EMA}_t = \alpha \cdot u_t + (1 - \alpha) \cdot \text{EMA}_{t-1}
\]

\texttt{burst\_streak} increments when $u_t > 1.5 \times \text{EMA}_{t-1}$.
Compression triggers when \texttt{burst\_streak} $\geq 3$, or when the projected
total exceeds $\texttt{max\_tokens} \times (1 - \texttt{safety\_margin})$.

\subsection{Atomic checkpoints}

\marginnote{\color{added} New in v2.1}

\begin{lstlisting}[caption={Atomic write — safe on crash}]
def _atomic_write(path: str, data: bytes) -> None:
    fd, tmp = tempfile.mkstemp(dir=directory, suffix=".tmp")
    with os.fdopen(fd, "wb") as fh:
        fh.write(data)
        fh.flush()
        os.fsync(fh.fileno())   # flush to physical storage
    os.replace(tmp, path)       # atomic on POSIX
\end{lstlisting}

The checkpoint file is always a valid JSON document, even if the process is
killed mid-write.

\subsection{Reproducibility bundle}

\marginnote{\color{added} New in v2.1}

\begin{lstlisting}[caption={Export and verify a run bundle}]
bundle_path = export_bundle(kernel, artifact)
# → /tmp/chaos_bundle/repro_YYYYMMDD_HHMMSS.zip
# Contains: artifact.json, history.jsonl, manifest.json (SHA-256)

ok, info = verify_bundle(bundle_path)
# → (True, {"ok": True, "errors": []})
\end{lstlisting}

A run that cannot be reproduced is not production-ready. -AoCH

%% ============================================================
\section{Comparison with Typical Agent Architectures}

\begin{center}
\begin{tabularx}{\linewidth}{@{} l X X @{}}
  \toprule
  \sffamily\bfseries Concern
    & \sffamily\bfseries Typical approach
    & \sffamily\bfseries Chaos Agent \\
  \midrule
  Event routing    & Scattered across components       & Single \texttt{emit()} \\
  Safety checks    & Duplicated per agent              & Central policy + AST audit \\
  Depth management & Hard limit / exception            & Graceful compression \\
  Token management & Unmonitored or hard cut           & EMA \texttt{TokenBudget} \\
  Cognitive layer  & Purely reactive                   & Explicit signal injection \\
  Output control   & Hard-coded logging                & Pluggable async callback \\
  Reproducibility  & Partial                           & Full Artifact + verified bundle \\
  Context size     & Unfiltered                        & \texttt{prepare\_context()} selector \\
  \bottomrule
\end{tabularx}
\end{center}

%% ============================================================
\section{v2.2 Additions}

\subsection{Two-layer sandbox}

\marginnote{\color{added} New in v2.2}

\begin{description}[style=nextline, leftmargin=0pt]
  \item[Layer 1 — subprocess (always active)]
    Clean environment (\texttt{env\_whitelist}), wall-clock \texttt{timeout},
    \texttt{RLIMIT\_CPU} and \texttt{RLIMIT\_NPROC} via \texttt{preexec\_fn}.
    Platform-portable, no root required.
  \item[Layer 2 — seccomp-bpf (optional)]
    Syscall allowlist applied in \texttt{preexec\_fn} if \texttt{libseccomp}
    is available and \texttt{policy.use\_seccomp=True}.
    Silently skipped on macOS and containers without the library —
    Layer 1 remains active regardless.
\end{description}

\begin{lstlisting}[caption={Sandbox entry point — one call, full pipeline}]
result = await guardian.run_code(code)
# 1. AST audit
# 2. symlink pre-check
# 3. subprocess + resource limits + optional seccomp
# 4. symlink post-check
\end{lstlisting}

\subsection{Symlink escape detection}

\marginnote{\color{added} New in v2.2}

Before and after execution, \texttt{\_check\_symlinks()} walks the working
directory and resolves every symlink via \texttt{os.path.realpath}. Any target
outside \texttt{allowed\_paths} aborts or invalidates the run.

\subsection{Extended AST audit coverage}

\marginnote{\color{changed} Extended in v2.2}

The audit visitor now catches additional bypass vectors:

\begin{lstlisting}[caption={New blocked patterns in v2.2}]
__import__("os")                   # direct dunder import call
ctypes.CDLL("libc.so")            # native library loading
os.execv("/bin/sh", [...])         # process replacement
obj.__builtins__                   # dunder dict access
obj.__globals__["exec"]("...")     # globals bypass
'__im' + 'port__'                  # concat obfuscation
\end{lstlisting}

\subsection{\texttt{forbid\_risky\_imports} — intentional design}

\marginnote{\color{changed} Documented in v2.2}

\texttt{forbid\_risky\_imports} is \textbf{off by default}. This is a deliberate
architectural choice, not an oversight:

\begin{quote}
  \itshape We block dangerous calls, not imports.
  \texttt{import os} is legitimate in countless programmes.
  \texttt{os.system(\ldots)} is not. Blocking the import would break
  legitimate utilities without meaningfully improving security
  — the call is what causes harm. — AoCH
\end{quote}

The rationale is documented in the module docstring, \texttt{policy.json},
the README, and here. Enable \texttt{forbid\_risky\_imports=True} only when
executing fully untrusted code where even importing is unacceptable.

\subsection{Tooling \& CI}

\marginnote{\color{added} New in v2.2}

\begin{center}
\begin{tabular}{@{} >{\sffamily\bfseries}l p{9cm} @{}}
  \toprule
  File & Purpose \\
  \midrule
  \texttt{tools/lint\_agent.py}   & CLI linter over \texttt{audit\_code()};
                                    exit 0/1/2; \texttt{--json} for CI. \\
  \texttt{tools/policy\_loader.py} & Load \texttt{GuardianPolicy} from
                                    \texttt{policy.json}; safe defaults on error. \\
  \texttt{policy.json}            & Declarative policy with \texttt{\_comment}
                                    keys explaining every decision. \\
  \texttt{tests/test\_agent.py}   & 21 unit tests: AST audit, sandbox utils,
                                    atomic write, TokenBudget. \\
  \texttt{.github/workflows/ci.yml} & Three jobs: AST Audit → Unit Tests
                                    → Bundle Integrity (sequential). \\
  \bottomrule
\end{tabular}
\end{center}

%% ============================================================
\section{Known Limitations \& Roadmap}

\begin{description}[style=nextline, leftmargin=0pt]
  \item[seccomp syscall allowlist is conservative]
    The allowlist may block legitimate syscalls on uncommon platforms.
    Extend the list in \texttt{\_make\_preexec()} as needed.

  \item[\texttt{RLIMIT\_NPROC} may fail in containers]
    Some containers lack user namespaces required for this limit.
    The failure is silently ignored; timeout (Layer 1) remains active.

  \item[\texttt{allowed\_paths} not enforced at OS level]
    Declared in policy, enforced by symlink checks — not by chroot/seccomp
    path filter.
    \textit{Stage 3: seccomp path whitelist or chroot integration.}

  \item[Reactive token management]
    \textit{Stage 4: predictive token budgeting.}

  \item[MDA integration not yet implemented]
    Pre-commit hooks, GPG bundle signing, CI audit log archiving.
    \textit{Stage 3 roadmap.}
\end{description}

%% ============================================================
\section{Quick Start}

\begin{lstlisting}[language=bash, caption={Installation and first run}]
git clone https://github.com/iVANIUS1/chaos-agent
cd chaos-agent
pip install httpx
export ANTHROPIC_API_KEY=sk-ant-...
python chaos_agent_v2_1.py
\end{lstlisting}

%% ============================================================
\section{License}

MIT — free to use, modify, and build upon. Attribution appreciated.

\vspace{8mm}
{\color{rule}\rule{\linewidth}{0.4pt}}
\vspace{4pt}

\noindent{\sffamily\small\color{note}
  Ivan Putna · Architect of Chaos ·
  \href{https://github.com/iVANIUS1}{GitHub} ·
  \href{https://www.linkedin.com/in/ivan-putna-3313a3381}{LinkedIn}
}

\end{document}
